\section*{Chapter 3: Spacetime, Curvature and Gravity}
\begin{description}[style=multiline,leftmargin=2em,labelwidth=1.5em]

% DOCX paragraph 2023
\item[1.]\hypertarget{note:03-01}{}For those interested in a more rigorous mathematical definition, spacetime is defined
as a \emph{variety} with events as elements. To this variety, it is possible to
associate a ``map'' of coordinates\index{coordinates} through which the invariant
distances between events can be measured. This map is described by a rank-two tensor,
called \emph{metric tensor}\index{tensor!metric} or simply
``\emph{metric}\index{metric|see{tensor, metric}}\emph{''}. Therefore, the distance between events is
measured by the scalar product, which is mediated by the metric, of the separation
between events and is a scalar invariant, i.e., the same in all coordinate systems.

% DOCX paragraph 2024
\item[2.]\hypertarget{note:03-02}{}To be precise, the \emph{m} in Equation \eqref{eq:03-02} indicates the so-called ``\emph{rest
mass''}, \index{mass!rest} i.e., the portion of mass that does not depend on the
object's state of motion. This rest mass has no value in the case of a photon, which
thus has a zero rest mass. A more general expression of Eq. \eqref{eq:03-03}, and that is valid for
all types of particles, is \emph{E}\textsuperscript{2} \emph{= m}\textsuperscript{2}
\emph{c}\textsuperscript{4} + \emph{p}\textsuperscript{2} \emph{c}\textsuperscript{2},
where \emph{p} is the \emph{linear} \emph{momentum} of the object. This expression is
also valid for a photon, which always has a non-zero momentum proportional to its
frequency.

% DOCX paragraph 2025
\item[3.]\hypertarget{note:03-03}{}For the interested readers, I can briefly add that the Einstein equations summarised
in Eq. \eqref{eq:03-03} are ten highly nonlinear, second-order, partial
differential\index{differential!equations} equations. They are written in
covariant\index{Einstein!covariant} form through the use of two rank-two tensors, namely the
Einstein tensor\index{Einstein!tensor} \emph{G\textsubscript{µν}} and the energy-momentum
tensor \emph{T\textsubscript{µν}}.

% DOCX paragraph 2026
\item[4.]\hypertarget{note:03-04}{}Note that the example of the sheet and the bowling ball has a logical flaw: it
describes gravity through the effect of the bowling ball, but the latter is
\emph{already} subject to a gravitational field. We will explore the question further in
this chapter, demonstrating how curvature can be measured even without the need for a
``sheet''. For the moment, it is sufficient to point out that this analogy is perfectly
valid. It is important to understand here that the presence of matter or energy produces
curvature.

% DOCX paragraph 2027
\item[5.]\hypertarget{note:03-05}{}If you want to do a little quiz, ask yourself: given two points A and B on the Earth's
surface, what path between A and B minimises or maximises the distance? The answer is
the great \emph{circle} that unites them: the circle passing through A, B and the centre
of the Earth. For example, if we imagine that the two points have the same longitude,
that path would follow the meridian that unites them and that passes through the North
and South Poles. But, of course, in most cases, going through the poles maximises the
distance~...

\end{description}
